% =====================================================================
% PISTES POUR DE FUTURS MEMOIRES D'ACTUARIAT
%
% ANNEXE EXIGEE PAR L'INSTITUT DES ACTUAIRES. Le paragraphe 5.1.h des
% Recommandations Jury applicables depuis le 1er avril 2026 impose deux
% annexes, dont celle-ci : << Pistes eventuelles pour de futurs memoires
% actuariels >>, qui << peuvent reprendre de facon non limitative des pistes
% d'enrichissement du memoire presente mais aussi des axes d'etudes
% abandonnees et n'ayant pas conduit a une impasse definitive >>, et qui
% << constituent un legs a l'entreprise ayant accueilli l'etudiant mais aussi
% un apport pour la communaute actuarielle >>.
%
% CREE LE 21 SEPTEMBRE 2026. Aucune piste inventee : chacune est un travail
% deja instruit dans le memoire et laisse ouvert pour une raison ecrite. Les
% deux seules grandeurs chiffrees reprises ici (66 % et 28,9 %) sont celles du
% chapitre de la donnee manquante, et les scripts cites sont les siens.
% =====================================================================

\chapter{Pistes pour de futurs mémoires d'actuariat}
\label{ann:pistes}

\begin{cle}
\textbf{Idée directrice.} Les pistes ci-dessous ne sont pas des suggestions générales. Chacune
correspond à un travail que ce mémoire a instruit, outillé, puis laissé ouvert pour une raison
qu'il énonce : une donnée qui n'existe pas encore, un protocole préparé et non exécuté, ou une
limite dont il démontre qu'elle est structurelle. Elles sont classées par ce qu'il resterait à
faire, du plus immédiatement exécutable au plus exploratoire.
\end{cle}

\section{Exécuter l'élicitation structurée préparée par ce mémoire}

L'annexe~\ref{ch:elicitation} décrit un protocole de Cooke complet, préparé et non exécuté : la
direction de la contagion n'étant pas identifiable sur la donnée disponible, le jugement d'experts
est la seule source qui puisse la fixer. Le questionnaire est rédigé, sa partie d'étalonnage
comporte dix questions de graine dont les valeurs vraies sont calculées et contrôlées, et la
convention d'orientation des liens y est fixée sans ambiguïté.

Ce qui manque est le panel. Un mémoire qui conduirait cette élicitation auprès de responsables de
la résilience opérationnelle disposerait donc d'un instrument prêt, et son apport propre porterait
sur ce que ce mémoire n'a pas pu faire : la constitution du panel, la mesure de la calibration des
experts sur les questions de graine, la pondération des avis qui en découle, et la confrontation
des poids obtenus aux directions que le corpus documentaire établit déjà.

\begin{attention}
\textbf{Une précaution à ne pas perdre.} Les valeurs vraies des questions de graine ne sont pas
publiées dans ce mémoire, et ne doivent pas l'être : les publier brûlerait l'instrument pour toute
élicitation ultérieure. Un futur mémoire qui reprendrait le protocole doit maintenir cette
séparation entre le questionnaire et son barème.
\end{attention}

\section{Documenter la dépendance entre gouvernance et gestion des incidents}

Le chapitre~\ref{ch:livrable} hiérarchise les données manquantes par la valeur d'information
qu'elles portent, et ce classement est très inégal. Fixer le seul sens de la dépendance entre la
gouvernance du risque TIC et la gestion des incidents lèverait $28{,}9\,\%$ de
l'ambiguïté de capital ; fixer celui du lien entre tests et prestataires tiers n'en lèverait rien.

La dépendance la plus précieuse est précisément celle que le corpus de rapports publics ne couvre
pas. La lecture documentaire a livré ce qu'elle pouvait livrer, et le reste exige de la donnée
d'entité : une chronologie interne d'incidents, horodatée à la survenance et catégorisée par
domaine de contrôle. Un mémoire conduit chez un assureur disposant de ce registre interne aurait
accès à ce que ce travail a dû traiter en bornes, et pourrait mesurer de combien la bande se
resserre réellement, là où ce mémoire ne peut en donner que la valeur théorique.

% SOURCES-SCRIPTS: 55

\section{Le second codage en aveugle du corpus de rapports post-incident}

La direction des dépendances est codée à partir de rapports post-incident officiels. Ce codage est
le fait d'un seul lecteur, et le mémoire le déclare comme une limite : un biais de narration
mesuré, non supposé.

La réception d'un second codage indépendant est outillée : le script d'accord inter-codeurs ne
fabrique rien en l'absence de donnée et attend le second jeu. Un futur travail qui ferait coder le
même corpus en aveugle par un second lecteur produirait la mesure d'accord qui manque, et
permettrait de dire si la direction retenue tient à la lecture ou au texte. C'est un apport modeste
en volume mais qui porte sur le maillon le plus contesté de la chaîne.

\section{Calibrer l'amplitude des liens, et non leur seul sens}

Le chapitre~\ref{ch:livrable} établit un résultat qui ferme une porte et en ouvre une autre : même
les dix directions connues, une bande de capital subsisterait, imputable à la seule amplitude des
coefficients. Connaître le sens d'un lien ne donne pas sa force. L'identification partielle n'est
donc pas un état provisoire que la donnée finirait par lever.

Toute la valeur d'information mesurée ici, jusqu'à $66\,\%$ de resserrement si l'on retenait le
sens documenté par les rapports publics, porte sur la direction. L'amplitude reste posée, et le mémoire
ne s'appuie que sur l'ordre des niveaux de propagation, le capital étant croissant en eux. Un
mémoire qui s'attaquerait à l'estimation des amplitudes, par exemple sur des séries d'incidents
d'une même entité où la sévérité transmise est observable, traiterait la moitié du problème que
celui-ci laisse intacte.

% SOURCES-SCRIPTS: 66

\section{La dérive d'échelle des pertes, mesurée et non modélisée}

Le backtest hors échantillon couvre correctement la forme de la queue et la loi de fréquence, mais
pas le niveau de la charge annuelle. L'écart s'explique par une dérive de l'échelle des pertes dans
le temps, que ce mémoire mesure et signale sans la modéliser : la calibration reste stationnaire.

Un mémoire portant sur la non-stationnarité de la sévérité cyber (inflation des coûts de
remédiation, évolution des montants de rançon, effet du durcissement des garanties) corrigerait
la seule composante du backtest que ce travail échoue à couvrir. Le diagnostic est posé et chiffré
ici, ce qui donne à un tel travail son point de départ et son critère de réussite.

\section{Le passage de l'échelle sectorielle à l'entité}

Le besoin de capital produit ici est ramené à l'échelle d'une entité au moyen de rapports publics
de solvabilité, et le mémoire publie la borne inférieure de validité de cette opération. C'est le
point où un lecteur d'entreprise demandera légitimement davantage.

Un mémoire conduit à l'intérieur d'une entité, avec son exposition réelle, son registre de
prestataires au sens du règlement d'exécution et son propre historique d'incidents, remplacerait
cet ancrage par une mesure directe. Il pourrait en outre traiter la question que ce travail ne fait
qu'ouvrir : comment un besoin de capital de type ORSA, sensible à l'état de conformité,
s'articule avec une Formule Standard qui ne comporte aucun module DORA.

\begin{cle}
\textbf{Ce qui se lègue avec ces pistes.} Les quatre premières sont exécutables avec ce qui existe
déjà au dossier : un protocole d'élicitation rédigé et étalonné, un classement chiffré des données
manquantes, un dispositif de second codage en attente de sa donnée, et un diagnostic de dérive posé
avec sa mesure. Les deux dernières demandent un accès à de la donnée d'entité, que ce mémoire n'a
pas eu. Aucune ne suppose de reprendre le modèle à zéro.
\end{cle}

% =====================================================================
